%! Roots and Polynomials — Unit 3, Lesson 3.6 — Guided Notes (ANSWER KEY)
% MIRROR of ../notes/main.tex: the vocabulary write lines -> \ansline, \blank -> \ans, the hook's
% \writelines{2} -> one \ansline, and the answer argument of every \pcell, \writespace and
% \labelbox filled. Nothing else differs. No teachernote here — teacher prose is in the plan.
% \ans{} and every \pcell answer are TEXT MODE: math inside them is wrapped in $...$.
\documentclass[10pt]{article}
\usepackage{precalculus-article}
\usepackage{precalculus-key}
\newcommand{\TallMath}[1]{$\displaystyle #1\rule[-1.4em]{0pt}{3.2em}$}

\begin{document}

\pageheader{Unit 3, Lesson 3.6}{Guided Notes --- Answer Key}

\vspace{0.12in}

\boxguard
\begin{vocabbox}
Fill in each term as we use it together.
\par\vspace{2pt}
\noindent\textbf{\textcolor{plum}{Zero and factor:}}\\[1pt]
\ansline{If $a$ is a zero of $p$, then $(x - a)$ is a factor of $p$ --- and the reverse.}\\[3pt]
\noindent\textbf{\textcolor{plum}{Multiplicity:}}\\[1pt]
\ansline{How many times a factor $(x - a)$ appears; the degree is the sum of the multiplicities.}\\[3pt]
\noindent\textbf{\textcolor{plum}{Factored form:}}\\[1pt]
\ansline{A leading coefficient times the factors, $a(x - r_1)(x - r_2)\cdots$; it shows the zeros.}\\[3pt]
\noindent\textbf{\textcolor{plum}{Standard form:}}\\[1pt]
\ansline{Terms in descending powers of $x$; it shows the degree, the leading term, the end behavior.}\\[3pt]
\noindent\textbf{\textcolor{plum}{Conjugate pair:}}\\[1pt]
\ansline{$a + bi$ and $a - bi$. With real coefficients, a non-real zero always brings its conjugate.}
\end{vocabbox}

\vspace{0.1in}

\boxguard
\begin{hookbox}
I am thinking of a polynomial. Its only zeros are $1$ and $4$. Write a polynomial that could
be mine. Is yours the only one that works? Name two more.
\ansline{$(x - 1)(x - 4)$ works --- but so do $3(x - 1)(x - 4)$, $-(x - 1)(x - 4)$, and $(x - 1)^2(x - 4)$. The zeros alone do not pin down one polynomial: the stretch and the multiplicities are still free.}
\end{hookbox}

\small
\begin{guidednotes}
% ── ROW 1: zeros -> factors -> standard form ─────────────────────────────────
\mainidea[Zeros to factors]{Build It} &
If $a$ is a zero of a polynomial, then $(x - a)$ is a factor --- so each zero hands you one
factor, written as $x$ \textbf{minus the zero}. A negative zero turns into a plus: the zero
$-2$ gives $(x + 2)$, because $-2 + 2 = 0$. Multiply the factors together, two at a time, to
reach \textbf{standard form}.

\vspace{4pt}
\begin{center}
\renewcommand{\arraystretch}{1.5}
\begin{tabular}{|>{\bfseries}l|c|c|c|c|c|}
\hline
Zero   & $3$         & $-2$         & $0$          & $-5$         & $a$ \\ \hline
Factor & $(x - 3)$   & \ans{$(x + 2)$} & \ans{$x$} & \ans{$(x + 5)$} & $(x - a)$ \\ \hline
\end{tabular}

\vspace{8pt}
\begin{tikzpicture}[>=stealth, font=\footnotesize,
    stage/.style={draw=plum, thick, fill=lilac, rounded corners=3pt, text width=2.45cm,
                  align=center, minimum height=1.25cm, inner sep=3pt}]
  \node[stage] (z)  at (0,0)    {\textbf{Zeros}\\ $-2,\ 1,\ 3$};
  \node[stage] (f)  at (3.15,0) {\textbf{One factor each}\\ $(x + 2)$, $(x - 1)$,\\ $(x - 3)$};
  \node[stage] (ff) at (6.3,0)  {\textbf{Factored form}\\ $(x + 2)(x - 1)(x - 3)$};
  \node[stage, fill=hookbg, draw=goldacc] (sf) at (9.45,0) {\textbf{Standard form}\\ multiply it out: problem 1};
  \draw[->, very thick, plum] (z) -- (f);
  \draw[->, very thick, plum] (f) -- (ff);
  \draw[->, very thick, plum] (ff) -- (sf);
  \node[font=\scriptsize\itshape, below] at (1.575,-0.62) {flip the sign};
  \node[font=\scriptsize\itshape, below] at (7.875,-0.62) {two at a time};
\end{tikzpicture}
\end{center}
\\
& \notesprompt{Leading coefficient 1. Write the factored form, then the standard form.}
\begin{probgrid}{|Y|Y|}
\pcell{1}{Zeros $-2$, $1$, and $3$.}{2.6cm}{$(x + 2)(x - 1)(x - 3)$\newline $= (x^2 + x - 2)(x - 3)$\newline $= x^3 - 2x^2 - 5x + 6$} &
\pcell{2}{Zeros $4$ and $-3$.}{2.6cm}{$(x - 4)(x + 3)$\newline $= x^2 + 3x - 4x - 12$\newline $= x^2 - x - 12$} \\ \hline
\end{probgrid}
\\
& \begin{probgrid}*{|Y|Y|}
\pcell{3}{Zeros $0$, $2$, and $-5$. What factor does the zero $0$ give?}{2.6cm}{The zero 0 gives the factor $x$.\newline $x(x - 2)(x + 5) = x(x^2 + 3x - 10)$\newline $= x^3 + 3x^2 - 10x$} &
\pcell{4}{Jordan wants zeros $-4$ and $1$ and writes $(x - 4)(x + 1)$. Test $x = -4$ in Jordan's product. Then fix it and give the standard form.}{2.6cm}{$(-4 - 4)(-4 + 1) = (-8)(-3) = 24 \ne 0$, so $-4$ is not a zero of Jordan's product.\newline Fixed: $(x + 4)(x - 1) = x^2 + 3x - 4$} \\ \hline
\end{probgrid}
\\ \hline
% ── ROW 2: multiplicity and the leading coefficient ──────────────────────────
\mainidea[Repeats and the stretch]{Fit the Point} &
A zero of \textbf{multiplicity} $k$ puts its factor in $k$ times, as $(x - a)^k$, and the
degree is the sum of the multiplicities. The zeros fix only a \emph{family}: every nonzero
number $a$ in front gives the same zeros and a different stretch. One more point on the graph
fixes $a$ --- substitute the point and solve.

\vspace{4pt}
\begin{center}
\begin{tikzpicture}[>=stealth]
  \node[font=\Large, anchor=east] (lhs) at (0,0) {$f(x) = $};
  \node[font=\Large, anchor=west, text=goldacc] (a) at (0.05,0) {$\boldsymbol{a}$};
  \node[font=\Large, anchor=west] (p1) at (1.2,0) {$(x - 2)^2$};
  \node[font=\Large, anchor=west] (p2) at (3.6,0) {$(x + 1)$};
  \node[font=\Large, anchor=west] (p3) at (5.8,0) {$(x - 3)$};
  \foreach \n/\t/\c in {a/{the stretch:\\ found from a point}/goldacc,
                        p1/{zero $2$,\\ multiplicity 2}/plum,
                        p2/{zero $-1$,\\ multiplicity 1}/plum,
                        p3/{zero $3$,\\ multiplicity 1}/plum}
    \draw[->, thick, \c] (\n.south) -- ++(0,-0.6) node[below, font=\scriptsize, align=center,
         text width=2.1cm, text=\c] {\t};
\end{tikzpicture}

\vspace{4pt}
Degree $= 2 + 1 + 1 = $ \labelbox{1.0cm}{$4$}
\end{center}

\stepnum{1}\ One factor per zero, with its multiplicity as the exponent.\par
\stepnum{2}\ Put $a$ in front.\par
\stepnum{3}\ Substitute the point and solve for $a$.
\\
& \notesprompt{Write each polynomial. Use the given point to find $a$.}
\begin{probgrid}{|Y|Y|}
\pcell{5}{Degree 4: zero $2$ with multiplicity 2, zeros $-1$ and $3$, and the point $(0, 24)$. Leave it in factored form.}{2.4cm}{$f(0) = a(-2)^2(1)(-3) = -12a = 24$, so $a = -2$.\newline $f(x) = -2(x - 2)^2(x + 1)(x - 3)$} &
\pcell{6}{Degree 3: zeros $-3$, $1$, and $2$, and the point $(0, 12)$. Factored form, then standard form.}{2.4cm}{$a(3)(-1)(-2) = 6a = 12$, so $a = 2$.\newline $2(x + 3)(x - 1)(x - 2) = 2(x^3 - 7x + 6)$\newline $= 2x^3 - 14x + 12$} \\ \hline
\end{probgrid}
\\
& \begin{probgrid}*{|Y|Y|}
\pcell{7}{Degree 3: zero $-1$ with multiplicity 2, zero $4$, and $y$-intercept $-8$.}{2.2cm}{$a(0 + 1)^2(0 - 4) = -4a = -8$, so $a = 2$.\newline $f(x) = 2(x + 1)^2(x - 4)$} &
\pcell{8}{Kai answers problem 7 with $(x + 1)^2(x - 4)$. Find the $y$-intercept of Kai's polynomial. What did he forget?}{2.2cm}{Kai's $y$-intercept is $(1)^2(-4) = -4$, not $-8$. His zeros are right, but he forgot $a$ --- the point is what fixes $a = 2$.} \\ \hline
\end{probgrid}
\\ \hline
% ── ROW 3: reading zeros, multiplicity, and a point off a graph ──────────────
\mainidea[Zeros you can see]{Off a Graph} &
Every $x$-intercept is a real zero. Where the graph \textbf{crosses} the $x$-axis the
multiplicity is odd (use 1); where it \textbf{touches} the axis and turns back --- tangent to
it --- the multiplicity is even (use 2). Then read one more labeled point, off the $x$-axis, to
find $a$.
\\
& \begin{center}
\begin{tikzpicture}[x=0.78cm, y=0.3cm, >=stealth]
  \draw[linegray, very thin, step=1] (-3,-5) grid (3,5);
  \draw[->, thick] (-3,0) -- (3.4,0) node[right] {\scriptsize $x$};
  \draw[->, thick] (0,-5) -- (0,5.6) node[above] {\scriptsize $y$};
  \foreach \x in {-2,-1,1,2} \node[below, font=\tiny] at (\x,0) {\x};
  \foreach \y in {-4,-2,2,4} \node[left, font=\tiny] at (0,\y) {\y};
  \draw[very thick, plum, domain=-2.45:2.1, samples=80, smooth]
       plot (\x, {-(\x+2)*(\x-1)*(\x-1)});
  \fill[plum] (-2,0) circle (2.2pt);
  \fill[plum] (1,0) circle (2.2pt);
  \fill[goldacc] (0,-2) circle (2.4pt) node[below right, font=\scriptsize, fill=white, inner sep=1pt] {$(0, -2)$};
  \node[font=\small\bfseries, text=plum] at (-2.1,4.3) {A};
\end{tikzpicture}
\hspace{0.6cm}
\begin{tikzpicture}[x=0.78cm, y=0.3cm, >=stealth]
  \draw[linegray, very thin, step=1] (-1,-5) grid (5,5);
  \draw[->, thick] (-1,0) -- (5.4,0) node[right] {\scriptsize $x$};
  \draw[->, thick] (0,-5) -- (0,5.6) node[above] {\scriptsize $y$};
  \foreach \x in {1,2,3,4} \node[below, font=\tiny] at (\x,0) {\x};
  \foreach \y in {-4,-2,2,4} \node[left, font=\tiny] at (0,\y) {\y};
  \draw[very thick, plum, domain=-0.65:4.35, samples=80, smooth]
       plot (\x, {0.5*\x*(\x-3)*(\x-3)});
  \fill[plum] (0,0) circle (2.2pt);
  \fill[plum] (3,0) circle (2.2pt);
  \fill[goldacc] (1,2) circle (2.4pt) node[above right, font=\scriptsize, fill=white, inner sep=1pt] {$(1, 2)$};
  \node[font=\small\bfseries, text=plum] at (4.5,-4.2) {B};
\end{tikzpicture}

\vspace{4pt}
Graph A at $x = 1$: \labelbox{1.8cm}{touches}\qquad Graph B at $x = 0$: \labelbox{1.8cm}{crosses}
\end{center}
\\
& \notesprompt{Each graph is a cubic. Write its equation in factored form.}
\begin{probgrid}{|Y|Y|}
\pcell{9}{Graph A.}{2.4cm}{Crosses at $-2$, touches at $1$: $f(x) = a(x + 2)(x - 1)^2$.\newline $f(0) = a(2)(1) = 2a = -2$, so $a = -1$.\newline $f(x) = -(x + 2)(x - 1)^2$} &
\pcell{10}{Graph B. Why is the $y$-intercept no help in finding $a$ this time?}{2.4cm}{Crosses at $0$, touches at $3$: $g(x) = ax(x - 3)^2$. The $y$-intercept is 0 for every $a$, so use $(1, 2)$: $4a = 2$, $a = \tfrac{1}{2}$.\newline $g(x) = \tfrac{1}{2}x(x - 3)^2$} \\ \hline
\end{probgrid}
\\ \hline
% ── ROW 4: THE CRUX — the forgotten conjugate ────────────────────────────────
\mainidea[Zeros you cannot see]{Conjugate Pairs} &
A polynomial of degree $n$ has exactly $n$ complex zeros, counting multiplicity --- but the
non-real ones never show up on a graph. If the coefficients are real and $a + bi$ is a zero,
then its \textbf{conjugate} $a - bi$ is a zero too. Multiply the pair \emph{first}: the $i$'s
cancel and leave one quadratic factor with real coefficients.
\\
& \begin{center}
\fcolorbox{plum}{lilac}{\begin{minipage}{0.78\linewidth}\centering\vspace{3pt}
\textbf{A conjugate pair becomes one real quadratic factor}
\[
\begin{aligned}
\bigl(x - (2 + i)\bigr)\bigl(x - (2 - i)\bigr) &= \bigl((x - 2) - i\bigr)\bigl((x - 2) + i\bigr) \\
  &= (x - 2)^2 - i^2 \\
  &= x^2 - 4x + 4 + 1 \\
  &= x^2 - 4x + 5
\end{aligned}
\]
\vspace{-2pt}
The $-i^2$ became $+1$ because $i^2 = $ \labelbox{1.0cm}{$-1$}
\vspace{3pt}
\end{minipage}}
\end{center}

\vspace{2pt}
\textbf{Real coefficients: non-real zeros come in conjugate pairs.} Each pair is one quadratic
factor, so it adds 2 to the degree --- never 1.

\vspace{2pt}
\textbf{In your own words:} why must $2 - i$ come along whenever $2 + i$ is a zero?
\writespace{1.2cm}{Without $2 - i$ the $i$'s would not cancel, and the polynomial would have a non-real coefficient. Real coefficients force the pair.}
\\
& \notesprompt{Least degree, real coefficients, leading coefficient 1.}
\begin{probgrid}{|Y|Y|}
\pcell{11}{Zeros $1$ and $2 + i$. Use the display; give the standard form.}{2.4cm}{$2 - i$ comes too: $(x - 1)(x^2 - 4x + 5)$\newline $= x^3 - 4x^2 + 5x - x^2 + 4x - 5$\newline $= x^3 - 5x^2 + 9x - 5$ \quad (degree 3)} &
\pcell{12}{Zero $3i$. Which other zero must come with it? Write the quadratic factor.}{2.4cm}{$-3i$.\newline $(x - 3i)(x + 3i) = x^2 - 9i^2$\newline $= x^2 + 9$} \\ \hline
\end{probgrid}
\\
& \begin{probgrid}*{|Y|Y|}
\pcell{13}{Zeros $3$ and $1 + 2i$. Ana writes $(x - 3)\bigl(x - (1 + 2i)\bigr)$; Ben writes $(x - 3)(x^2 - 2x + 5)$. What is the degree of each? Which one has real coefficients?}{2.6cm}{Ana: degree 2, and it expands to $x^2 - (4 + 2i)x + (3 + 6i)$ --- non-real coefficients. She forgot the conjugate $1 - 2i$. Ben: degree 3, real coefficients. Ben is right.} &
\pcell{14}{Check that $(1 + 2i)(1 - 2i)$ gives Ben's 5. Then write Ben's polynomial in standard form.}{2.6cm}{$(1 + 2i)(1 - 2i) = 1 - 4i^2 = 5$\newline $(x - 3)(x^2 - 2x + 5)$\newline $= x^3 - 2x^2 + 5x - 3x^2 + 6x - 15$\newline $= x^3 - 5x^2 + 11x - 15$} \\ \hline
\end{probgrid}
\\ \hline
% ── LAST ROW: YOU DO — alone, while the teacher circulates (15 min) ──────────
\mainidea[You do, alone]{Build Four} &
Four polynomials to build from their zeros. Run the same three checks every time.
\\
& \begin{center}
\begin{tikzpicture}[>=stealth, font=\footnotesize,
    check/.style={draw=plum, thick, fill=lilac, rounded corners=3pt, text width=3.3cm,
                  align=center, minimum height=1.0cm, inner sep=3pt}]
  \node[check] (c1) at (0,0)   {\textbf{Each zero $\to$ a factor}\\ $x$ minus the zero: watch the sign};
  \node[check] (c2) at (4.0,0) {\textbf{Non-real zero?}\\ bring its conjugate; multiply the pair first};
  \node[check] (c3) at (8.0,0) {\textbf{A point given?}\\ put $a$ in front; substitute; solve};
  \draw[->, very thick, plum] (c1) -- (c2);
  \draw[->, very thick, plum] (c2) -- (c3);
\end{tikzpicture}
\end{center}
\\
& \notesprompt{On your own. Show your work.}
\begin{probgrid}{|Y|Y|}
\pcell{15}{Zeros $-4$, $1$, and $2$; leading coefficient 1. Factored form, then standard form.}{2.6cm}{$(x + 4)(x - 1)(x - 2)$\newline $= (x^2 + 3x - 4)(x - 2)$\newline $= x^3 + x^2 - 10x + 8$} &
\pcell{16}{Degree 3: zero $1$ with multiplicity 2, zero $-3$, and the point $(0, 6)$. Factored form.}{2.6cm}{$a(0 - 1)^2(0 + 3) = 3a = 6$, so $a = 2$.\newline $f(x) = 2(x - 1)^2(x + 3)$} \\ \hline
\end{probgrid}
\\
& \begin{probgrid}*{|Y|Y|}
\pcell{17}{Least degree, real coefficients, leading coefficient 1, zeros $-2$ and $3i$. State the degree, then give the standard form.}{2.8cm}{$3i$ brings $-3i$, so the degree is 3.\newline $(x + 2)(x^2 + 9)$\newline $= x^3 + 2x^2 + 9x + 18$} &
\pcell{18}{A degree-4 polynomial with real coefficients and leading coefficient 1 has zeros $2$, $-2$, and $1 - i$. Find its fourth zero, then write it in standard form.}{2.8cm}{Fourth zero: $1 + i$.\newline $(x^2 - 4)(x^2 - 2x + 2)$\newline $= x^4 - 2x^3 + 2x^2 - 4x^2 + 8x - 8$\newline $= x^4 - 2x^3 - 2x^2 + 8x - 8$} \\ \hline
\end{probgrid}
\\ \hline
\end{guidednotes}

% ── THE NOTES END AT THE YOU DO ROW; AP Practice follows; homework is DeltaMath ─

\end{document}
